% PDF page 67

\begin{center}
\rule{0.8\textwidth}{0.4pt}\\[0.6em]
{\large\bfseries APPENDIX}\\[0.3em]
{\Large\bfseries CLASSIFYING 1-MANIFOLDS}\\[0.6em]
\rule{0.8\textwidth}{0.4pt}
\end{center}

\noindent\textbf{WE WILL} prove the following result, which has been assumed in the text. A brief discussion of the classification problem for higher dimensional manifolds will also be given.

\noindent\textbf{Theorem.} \textit{Any smooth, connected $1$-dimensional manifold is diffeomorphic either to the circle $S^1$ or to some interval of real numbers.}

(An interval is a connected subset of $\mathbb{R}$ which is not a point. It may be finite or infinite; closed, open, or half-open.)

Since any interval is diffeomorphic\footnote{For example, use a diffeomorphism of the form $f(t) = a \tanh(t) + b$.} either to $[0, 1]$, $(0, 1]$, or $(0, 1)$, it follows that there are only four distinct connected $1$-manifolds.

The proof will make use of the concept of ``arc-length.'' Let $I$ be an interval.

\textsc{Definition.} A map $f : I \to M$ is a \emph{parametrization by arc-length} if $f$ maps $I$ diffeomorphically onto an open subset$\dagger$ of $M$, and if the ``velocity vector'' $\dd f_s(1) \in TM_{f(s)}$ has unit length, for each $s \in I$.\footnote{$\dagger$ Thus $I$ can have boundary points only if $M$ has boundary points.}

Any given local parametrization $I' \to M$ can be transformed into a parametrization by arc-length by a straightforward change of variables.

\noindent\textbf{Lemma.} \textit{Let $f : I \to M$ and $g : J \to M$ be parametrizations by arc-length. Then $f(I) \cap g(J)$ has at most two components. If it has only one component, then $f$ can be extended to a parametrization by arc-length of the union $f(I) \cup g(J)$. If it has two components, then $M$ must be diffeomorphic to $S^1$.}


% PDF page 68


\noindent\textit{Proof.}
Clearly $g^{-1} \circ f$ maps some relatively open subset of $I$ diffeomorphically onto a relatively open subset of $J$. Furthermore the derivative of $g^{-1} \circ f$ is equal to $\pm 1$ everywhere.

Consider the graph $\Gamma \subset I \times J$, consisting of all $(s, t)$ with $f(s) = g(t)$. Then $\Gamma$ is a closed subset of $I \times J$ made up of line segments of slope $\pm 1$. Since $\Gamma$ is closed and $g^{-1} \circ f$ is locally a diffeomorphism, these line segments cannot end in the interior of $I \times J$, but must extend to the boundary. Since $g^{-1} \circ f$ is one-one and single valued, there can be at most one of these segments ending on each of the four edges of the rectangle $I \times J$. Hence $\Gamma$ has at most two components. (See Figure~\ref{fig:milnor-19}.) Furthermore, if there are two components, the two must have the same slope.

\begin{figure}[htbp]
\centering
\setcounter{figure}{18}%
\IfFileExists{output/figures/fig19.tex}{% Figure 19: Three of the possibilities for Gamma
\begin{tikzpicture}[line cap=round,line join=round,>=Stealth,scale=0.9]
  % left-hand rectangle: two components of slope +1
  \draw[thick] (0,0) rectangle (2.6,2.6);
  % y-axis ticks: gamma, delta, alpha, beta (bottom to top)
  \foreach \y/\lab in {0.3/\gamma,0.9/\delta,1.7/\alpha,2.3/\beta} {
    \draw (-0.08,\y) -- (0.08,\y);
    \node[left] at (-0.14,\y) {\footnotesize $\lab$};
  }
  % x-axis ticks: a, b, c, d
  \foreach \x/\lab in {0.3/a,0.9/b,1.7/c,2.3/d} {
    \draw (\x,-0.08) -- (\x,0.08);
    \node[below] at (\x,-0.14) {\footnotesize $\lab$};
  }
  \draw[very thick] (0.3,1.7) -- (0.9,2.3);
  \draw[very thick] (1.7,0.3) -- (2.3,0.9);
  \node at (3.3,1.3) {or};
  % middle: a single component crossing a narrow rectangle
  \draw[thick] (3.9,0.55) rectangle (5.0,2.05);
  \draw[very thick] (3.9,0.7) -- (5.0,1.8);
  \node at (5.6,1.3) {or};
  % right: a long rectangle, one component, and further possibilities
  \draw[thick] (6.2,0.85) rectangle (8.9,1.75);
  \draw[very thick] (7.85,0.85) -- (8.75,1.75);
  \node at (9.35,1.3) {$\cdots$};
\end{tikzpicture}
}{% Figure 19: Three of the possibilities for Gamma
\begin{tikzpicture}[line cap=round,line join=round,>=Stealth,scale=0.9]
  % left-hand rectangle: two components of slope +1
  \draw[thick] (0,0) rectangle (2.6,2.6);
  % y-axis ticks: gamma, delta, alpha, beta (bottom to top)
  \foreach \y/\lab in {0.3/\gamma,0.9/\delta,1.7/\alpha,2.3/\beta} {
    \draw (-0.08,\y) -- (0.08,\y);
    \node[left] at (-0.14,\y) {\footnotesize $\lab$};
  }
  % x-axis ticks: a, b, c, d
  \foreach \x/\lab in {0.3/a,0.9/b,1.7/c,2.3/d} {
    \draw (\x,-0.08) -- (\x,0.08);
    \node[below] at (\x,-0.14) {\footnotesize $\lab$};
  }
  \draw[very thick] (0.3,1.7) -- (0.9,2.3);
  \draw[very thick] (1.7,0.3) -- (2.3,0.9);
  \node at (3.3,1.3) {or};
  % middle: a single component crossing a narrow rectangle
  \draw[thick] (3.9,0.55) rectangle (5.0,2.05);
  \draw[very thick] (3.9,0.7) -- (5.0,1.8);
  \node at (5.6,1.3) {or};
  % right: a long rectangle, one component, and further possibilities
  \draw[thick] (6.2,0.85) rectangle (8.9,1.75);
  \draw[very thick] (7.85,0.85) -- (8.75,1.75);
  \node at (9.35,1.3) {$\cdots$};
\end{tikzpicture}
}
\caption{Three of the possibilities for $\Gamma$.}\label{fig:milnor-19}
\end{figure}

If $\Gamma$ is connected, then $g^{-1} \circ f$ extends to a linear map $L : \mathbb{R} \to \mathbb{R}$. Now $f$ and $g \circ L$ piece together to yield the required extension
\[
F : I \cup L^{-1}(J) \to f(I) \cup g(J).
\]
If $\Gamma$ has two components, with slope say $+1$, they must be arranged as in the left-hand rectangle of Figure~\ref{fig:milnor-19}. Translating the interval $J = (\gamma, \beta)$ if necessary, we may assume that $\gamma = c$ and $\delta = d$, so that
\[
a < b \le c < d \le \alpha < \beta.
\]
Now setting $\theta = 2\pi t / (\alpha - a)$, the required diffeomorphism
\[
h : S^1 \to M
\]
is defined by the formula
\[
\begin{aligned}
h(\cos \theta, \sin \theta) &= f(t) && \text{for } a < t < d, \\
&= g(t) && \text{for } c < t < \beta.
\end{aligned}
\]


% PDF page 69

The image $h(S^1)$, being compact and open in $M$, must be the entire manifold $M$. This proves the lemma.

\noindent\textit{Proof of Classification Theorem.}
Any parametrization by arc-length can be extended to one
\[
f : I \to M
\]
which is maximal in the sense that $f$ cannot be extended over any larger interval as a parametrization by arc-length: it is only necessary to extend $f$ as far as possible to the left and then as far as possible to the right.

If $M$ is not diffeomorphic to $S^1$, we will prove that $f$ is onto, and hence is a diffeomorphism. For if the open set $f(I)$ were not all of $M$, there would be a limit point $x$ of $f(I)$ in $M - f(I)$. Parametrizing a neighborhood of $x$ by arc-length and applying the lemma, we would see that $f$ can be extended over a larger interval. This contradicts the assumption that $f$ is maximal and hence completes the proof.

\noindent\textbf{REMARKS.} For manifolds of higher dimension the classification problem becomes quite formidable. For $2$-dimensional manifolds, a thorough exposition has been given by Ker\'ekj\'art\'o [17]. The study of $3$-dimensional manifolds is very much a topic of current research. (See Papakyriakopoulos [26].) For compact manifolds of dimension $\ge 4$ the classification problem is actually unsolvable.\footnote{See Markov [19]; and also a forthcoming paper by Boone, Haken, and Po\'enaru in \emph{Fundamenta Mathematicae}.} But for high dimensional simply connected manifolds there has been much progress in recent years, as exemplified by the work of Smale [31] and Wall [37].

